\chapter{Receiving bytes without inventing boundaries}
\chapterintro{A correct server must tolerate a header arriving in any number of pieces. We replace the single-read assumption with an explicit buffer, delimiter search, and limit.}

\begin{objectives}
\begin{itemize}
  \item Accumulate partial reads until an HTTP header delimiter is complete.
  \item Treat connection close before completion as a protocol error.
  \item Enforce a cumulative byte limit rather than relying on receive-buffer size.
  \item Identify the bytes that belong after the header section.
\end{itemize}
\end{objectives}

\section{The tempting incorrect version}

A common first implementation calls \code{recv(4096)} and parses the result as
one request. The number 4096 is a maximum returned by that call, not a promise
that the peer's complete message fits or arrives together. A valid request can
be split after any byte.

Another incorrect version loops until a receive returns fewer than 4096 bytes.
A short read does not mark a message boundary. It only describes what was
available at that moment. The client can send more later.

HTTP supplies the header boundary: \code{b"\\r\\n\\r\\n"}. Therefore the
algorithm accumulates until it finds that exact sequence.

\section{The buffered snapshot}

\codefile{../code/stages/02_buffered_request/server.py}

The \code{bytearray} is mutable, so extending it does not create a new bytes
object on every receive. Each loop iteration has four possible outcomes:

\begin{enumerate}
  \item the delimiter is already present and the loop stops;
  \item more bytes arrive and are appended;
  \item an empty result reports end-of-stream before completion;
  \item the cumulative size exceeds the configured limit.
\end{enumerate}

After finding the delimiter, the function returns through its end. Notice a
design flaw that is acceptable only at this stage: bytes after the delimiter
are discarded. Those bytes might be a request body or the next request. Chapter
10 introduces a persistent incremental parser whose buffer survives and consumes
exactly one framed message.

\begin{pitfall}
The receive chunk size and protocol limit solve different problems. A 4 KiB
chunk is an I/O choice. A 32 KiB header limit is a policy applied to the total
accumulated header. Increasing one does not correctly implement the other.
\end{pitfall}

\section{End-of-stream is data about the stream}

When \code{recv} returns \code{b""}, the peer has performed an orderly close in
that direction. If the parser has no buffered bytes, the connection ended
between requests. If it has the beginning of a request, the message is incomplete
and should not be dispatched.

A reset or other socket error follows a different exception path. In both cases
the application must stop treating the stream as a source of a valid request.
There is no invented ``best effort'' completion.

\section{Delimiter search and strict CRLF}

Searching for the correct delimiter finds a candidate header end. A strict parser
must still reject bare line feeds or stray carriage returns inside the start and
field lines. Merely finding one correct final delimiter does not make every
earlier line correct.

Why be strict? Tolerant recipients often tolerate different things. If a reverse
proxy treats a bare newline as a boundary while the origin does not, their stream
interpretations diverge. Strictness reduces the space of ambiguous inputs.

\section{Timeouts belong beside byte limits}

A size limit stops a peer that sends too much. It does not stop a peer that sends
one permitted byte every minute. A read timeout limits how long a worker waits
for progress or for the next request. Later, each connected socket receives a
five-second timeout and incomplete timed-out requests receive 408 when practical.

Timeout values express a service policy. Too short rejects legitimate slow
clients; too long preserves resources for abusive or failed clients. Production
systems coordinate timeouts across proxies, application servers, and upstream
dependencies.

\begin{tryit}
Use a tiny client or interactive TCP tool to type a request slowly. Observe that
stage 2 waits because it has no timeout. Then mentally apply the final policy:
after five seconds without completion, close the connection and release the
worker.
\end{tryit}

\section{Design the next buffer}

The final parser needs an operation with different ownership:

\begin{itemize}
  \item \code{feed(data)} appends newly received bytes and applies early limits;
  \item \code{pop\_request()} returns one complete request or \code{None};
  \item parsed bytes are removed, while later bytes remain;
  \item parser state belongs to one connection, not one receive call.
\end{itemize}

This API separates I/O from syntax. Unit tests can feed fragments without opening
a real socket. The server loop can ask for another chunk only when the parser
reports that the current buffer is incomplete.

\summarytitle

Receive operations expose stream availability, not HTTP boundaries. A correct
reader accumulates bytes, searches for protocol delimiters, handles early close,
and enforces cumulative limits and timeouts. It must preserve bytes beyond one
message for later parsing.

\reviewtitle
\begin{enumerate}
  \item Why is a short receive not an end-of-message signal?
  \item What is wrong with discarding bytes after the header delimiter?
  \item Which attacks are addressed by byte limits and which by timeouts?
  \item Why should parser state belong to a connection?
\end{enumerate}

\begin{exercise}[Fragment table]
Take a small GET request and divide it at every possible byte boundary. For each
division, show the buffer before and after the second fragment. The parsing
result must be identical for every division.
\end{exercise}

\begin{exercise}[Limit placement]
Explain where to enforce request-line, total-header, field-count, and body limits.
For each, state the earliest moment the server has enough information to reject
the request without waiting for more bytes.
\end{exercise}

